Safety-critical work in chemical process plants is rarely the job of a single operator. The board operator at the console and the field operator on the unit coordinate through a continuous verbal channel over a shared mental model of the plant, and incident reviews across the petrochemical industry are consistent in implicating the link between operators rather than failure of either operator alone. The coordination axis between the two operators, in other words, carries an outsized share of the risk in process safety, and it is the axis that single-operator monitoring instruments cannot see. This thesis is a methodological precursor to operator-pair monitoring for chemical-engineering control rooms. It develops, on a laboratory analogue of the board-field coordination problem (an instructor-and-follower map task in which one operator verbally guides the other along a route under a shared spatial referent), the sensing apparatus, feature pipeline, and predictive model that the chemical-engineering deployment will require. Five operator-side modalities are recorded synchronously from both members of the pair: cardiac, gaze, audio, post-task survey, and per-trial dialogue annotation. Two screen-mounted eye-trackers operate in parallel (one per operator), and both gaze streams are projected into a shared map coordinate frame so that joint visual-attention measures across the role-asymmetric operators become definable. On forty pairs, the multimodal predictor reaches a cross-pair AUC of $0.76$ for trial-level coordination-failure prediction on previously unseen pairs, against $0.65$ for the best single modality. Three within-pair-normalised workload composites characterise per-operator and pair-level cognitive load, and a four-mode physiology-derived failure-mode taxonomy (Director-Overloaded, Matcher-Disengaged, Director-Disengaged, and Calm-Decoupled) decomposes the predictor's failure surface into operator-side mechanisms that map directly onto the overload and vigilance-decrement constructs of the abnormal-situation-management literature. Joint cardiac coupling separates failure from success at Cohen's $d \approx 0.7$ across three independent coupling features, indicating that even the lowest-friction operator-side sensor in a control-room deployment, the wrist smartwatch on both operators, would already carry meaningful coordination-state information. The contribution of this thesis is therefore three-fold: a multimodal sensing apparatus that resolves the role-asymmetric dual eye-tracking problem; an empirical demonstration that operator-pair coordination failure is predictable at trial-level resolution from a heterogeneous reference cohort; and a failure-mode taxonomy and workload-composite framework that maps onto the constructs already used in chemical-engineering process-safety monitoring, against which a confirmatory study with trained plant operators is the natural next step. Apparatus design, feature pipelines, and analysis code are released alongside the thesis to support that confirmatory study and subsequent operational deployment, at \href{https://github.com/valiant-disciple/map-task}{\texttt{github.com/valiant-disciple/map-task}} (experiment UI and data-collection stack) and \href{https://github.com/valiant-disciple/operator-pair-analysis}{\texttt{github.com/valiant-disciple/operator-pair-analysis}} (analysis pipeline and thesis source).
